\section{One service, pushed}
\label{sec:ch05-one-service-pushed}
\index{subscriptions|(}

A customer writes a review and somebody watching that product's page should see
it. That is the whole requirement, and meeting it needs the one operation type
Mosaic has not used: a query and a mutation each produce one response and close,
where a subscription opens a stream and delivers one result per event for as
long as the client stays connected~\autocite{chillicream2026subscriptions}.

Mosaic gets exactly one, and the scope is deliberate. It is a single service pushing to
clients connected to that service, over a pub/sub that lives in the process. The
federated version, where a router holds the client connection and a subgraph
publishes the event, is chapter~\ref{ch:real-time-in-a-federated-world}'s, and
almost nothing here survives that trip unchanged. What does survive is the
shape of the field and the discipline about the topic.

\begin{graphqlsdl}
type Subscription {
  onReviewAdded(productId: ID!): Review!
}
\end{graphqlsdl}

\subsection*{Two methods, and only one of them runs per event}

\index{subscriptions!subscribe resolver}%
A subscription field is two functions and the split matters more than it looks.

\begin{minted}[fontsize=\footnotesize]{csharp}
public static ValueTask<ISourceStream<Review>> SubscribeToReviewsAsync(
    [ID<Product>] Guid productId,
    ITopicEventReceiver receiver,
    CancellationToken cancellationToken)
    => receiver.SubscribeAsync<Review>(
        ReviewTopics.ReviewAdded(productId),
        cancellationToken);

[Subscribe(With = nameof(SubscribeToReviewsAsync))]
public static Review OnReviewAdded(
    [ID<Product>] Guid productId,
    [EventMessage] Review review)
    => review;
\end{minted}

The first runs once, when the client subscribes, and its job is to produce the
stream. The second runs once per event, receives the published message through
\code{[EventMessage]}, and returns what the client asked for. Returning the
message unchanged is the common case; the hook exists so a payload can be
reshaped or an event dropped before it reaches anybody.

\index{subscriptions!topics}%
The documented way to filter by product is a placeholder in a \code{[Topic]}
attribute, which interpolates an argument value into the topic
name~\autocite{chillicream2026subscriptions}. It reads better than the above and
I did not use it, for one reason: a topic is a contract between two pieces of
code that never call each other, and with a placeholder the exact string a
\code{Guid} formats into is something you establish by running it. So the topic
is a function, and both ends call it:

\begin{minted}{csharp}
public static string ReviewAdded(Guid productId) => $"reviewAdded:{productId}";
\end{minted}

That is worth more than it costs. The publisher and the subscriber cannot drift,
and the compiler is the thing keeping them together rather than a convention.

\subsection*{Publishing, and when}

\index{subscriptions!ITopicEventSender}%
The mutation from section~\ref{sec:ch05-errors-that-are-data} ends like this:

\begin{minted}[fontsize=\footnotesize]{csharp}
var review = await reviews.SubmitReviewAsync(
    productId, customerId, rating, body, cancellationToken);

await sender.SendAsync(
    ReviewTopics.ReviewAdded(productId),
    review,
    cancellationToken);

return review;
\end{minted}

\code{ITopicEventSender} is injected like any other service, and the abstraction
is over the provider, so the same line publishes to an in-memory bus, to Redis
or to NATS depending only on registration~\autocite{chillicream2026subscriptions}.

The position of that call is the decision. It is after the write has committed,
not before. A subscriber told about a review that then fails to save has been
lied to, and there is no message that takes it back. Publishing after the commit
means a crash between the two loses an event, and losing an event is recoverable
in a way that inventing one is not: a client that missed a notification can
refetch, and a client that acted on a notification about a row that does not
exist cannot.

\index{subscriptions!at-most-once}%
Which is the honest characterisation of what Mosaic delivers. At most once, no
durability, no replay. That is a real limitation and not a rounding error, and
it is the same limitation almost every in-process pub/sub has.

\subsection*{The registration, and the line that fails late}

\index{subscriptions!in-memory provider}%
Two lines of registration hold all of this up. One chooses the bus:

\begin{minted}{csharp}
builder.AddGraphQL()
    .AddMosaic()
    .AddInMemorySubscriptions()
\end{minted}

The other goes in the request pipeline:

\begin{minted}{csharp}
app.UseWebSockets();
\end{minted}

One warning about the second. Server-sent events already work
through \code{MapGraphQL}, which is the transport everything below is measured
over. The graphql-ws protocol over WebSocket is the one that needs
\code{UseWebSockets}~\autocite{chillicream2026subscriptions}, and the line is
here on the documentation's word rather than on mine: I did not exercise the
WebSocket transport, so I can tell you the registration is required and not what
the failure looks like without it. What is worth noticing either way is that
only one of the two transports depends on a line you can forget, and the schema
will build and the service will start regardless.

The provider is a package of its own,
\code{HotChocolate.Subscriptions.InMemory}, rather than something already in
\code{HotChocolate.AspNetCore}. Redis, NATS, RabbitMQ and PostgreSQL providers
sit beside it, and switching is a registration change with no effect on the
field or the publisher.

\subsection*{Watching it happen}

\index{subscriptions!server-sent events}%
Server-sent events are the easier of the two to try, because they are an
ordinary HTTP POST that does not end. Open one:

\begin{minted}[fontsize=\footnotesize]{text}
curl -N -X POST http://localhost:5100/graphql \
  -H 'Content-Type: application/json' \
  -H 'Accept: text/event-stream' \
  -d '{"query":"subscription { onReviewAdded(productId: \"...\")
                { id rating body author { displayName } } }"}'
\end{minted}

Send the mutation from another terminal. It answers first, as a mutation does:

\begin{minted}[fontsize=\footnotesize]{json}
{"data":{"submitReview":{"review":{"id":"UmV2aWV3Og7knwGk/wp1qg9QZb4xFzc="},
          "errors":null}}}
\end{minted}

And the open connection produces this:

\begin{minted}[fontsize=\footnotesize]{text}
event: next
data: {"data":{"onReviewAdded":{"id":"UmV2aWV3Og7knwGk/wp1qg9QZb4xFzc=",
       "rating":5,"body":"Boils fast, pours clean.",
       "author":{"displayName":"Sofia Ferrari"}}}}
\end{minted}

\index{subscriptions!execution per event}%
Two things in that payload are worth more than the fact that it arrived.

The identifier is the same one the mutation's own payload returned, which is
what you would hope and is not automatic: it is the same object, published
rather than re-read.

And \code{author} is populated. Nothing published a customer. The event carried
a \code{Review}, and \code{displayName} came back because the subscription ran
the selection set against that message, which meant running
\code{Review.author}, which meant a DataLoader executing inside the
subscription's own execution. A subscription is not a notification with a
payload attached. It is a query executed per event, with the event as its root
value, and every resolver behind every field runs each time. On a chatty topic
with a wide selection set that is a per-event cost that nobody budgeted for.

\subsection*{What breaks the moment there are two of them}

\index{subscriptions!single instance}%
Everything above works because the publisher and the subscriber are the same
process. Run two copies of Mosaic behind a load balancer and it stops: a client
connected to the first instance hears nothing about a review submitted through
the second, because the bus is a data structure in memory and there are now two
of them~\autocite{chillicream2026subscriptions}.

The fix is a registration change to a provider with a broker behind it, and the
reason it is not in this chapter is that swapping the provider is the easy part.
What follows it is not: connection affinity, what a client is owed when its
connection drops and reconnects somewhere else, whether events that happened
during the gap are replayed or lost, and, once a router is in front, which
process is holding the socket at all. Those are
chapter~\ref{ch:real-time-in-a-federated-world}'s, and they are most of the
subject.

Say the limit out loud rather than leaving it implied. In-memory subscriptions
are correct for exactly one instance, and one instance is a decision you have
made whether or not you noticed making it.

\index{subscriptions!not gated}%
One more admission, and it is about this book's own machinery rather than about
HotChocolate. Everything else in this chapter is asserted by
\code{scripts/verify.ps1} or by the Postman collection, and the subscription is
not. Newman speaks HTTP request and response, a subscription is a connection
that stays open, and I did not build a second harness for one field. So the two
transcripts above were produced by hand with \code{curl} and are the only
evidence in this chapter that nothing re-checks. If a later chapter breaks
\code{onReviewAdded}, the gate will pass. That is a real hole, it belongs to
chapter~\ref{ch:the-postman-federation-workbook}, where subscription testing is
the subject, and it is better written down than discovered.

Which is the schema this chapter set out to build: identifiers a client can
hold, a field with a bound on it, an answer for no, and a way to say something
happened. Every one of those is a promise that has to survive the graph being
split into six services, and chapter~\ref{ch:the-federation-model} starts
finding out which of them do.
\index{subscriptions|)}
